\scanpage{75}
\mspage{1}
% archive foliation: 74

Soit $E$ une courbe elliptique --- on travaille sur un corps alg. clos $k$ pour simplifier le propos, mais la seule hypothèse essentielle est que la base soit à car. résiduelles $\ne2$. On désigne par \deleted{$\unsure{\Sigma|E}$}
\[
(1)\qquad J={}_2E(k)^*
\]
l'ens des pts de $E$ d'ordre $2$, de sorte que l'on a
\[
(2)\qquad E/\{\pm\operatorname{id}_E\}\simeq\Sigma_J
\qquad\text{\deleted{\unsure{sphère}} \inserted{droite projective} $J$-marquée.}
\]
\[
(3)\qquad\Sigma_J\simeq\mathbb P(V_J(k))
\qquad 0\mathbin{-}V_J(k)\mathbin{-}k^J\xrightarrow{\text{somme}}k\longrightarrow0
\]
La donnée de $E$ est équivalente à celle de $J$ i.e. de $\Sigma_J$ muni du diviseur étale $D_J$ de degré $3$, plus le point
\[
(4)\qquad t\in\Sigma_J^*=\Sigma_J\setminus D_J
\]
image de l'élément $0_E$, plus une droite $T$ \inserted{sur $k$} (savoir $T_{E,0}$), munie d'un isomorphisme
\[
(5)\qquad T^{\otimes2}\xrightarrow{\sim}T_{\Sigma,t}
\]
(\deleted{isomorphisme} \inserted{équivalence} de catégories-groupoïdes \ldots valable en fait sur base $S$ quelconque à car. rés. $\ne2$, et compatible aux changements de base \ldots).

1 Notons qu'on a un isomorphisme canonique
\[
(6)\qquad T_\Sigma\simeq\underline O_\Sigma(2)(\underline\omega)
\]
où $\underline\omega=\underline\omega(J)$ est l'ens. des deux ordres circulaires sur $J$, et où $L(\underline\omega)$ (pour un objet \unsure{$M$|$L$} \unsure{que|en} faisceaux abéliens etc) désigne $L\wedge_{\{\pm1\}}\underline\omega$.
